% Chương 5 — Kết luận. Nạp bằng % Chương 5 — Kết luận. Nạp bằng % Chương 5 — Kết luận. Nạp bằng % Chương 5 — Kết luận. Nạp bằng \input{chapter5}. Chỉ dùng các gói đã khai báo.
\chapter{KẾT LUẬN VÀ HƯỚNG PHÁT TRIỂN}
\label{ch:ch5-ketluan}

\section{Kết quả đạt được}

Đồ án đã hoàn thành một hệ thống truy xuất văn bản tiếng Việt chạy hoàn toàn ngoại
tuyến, gồm ba tầng trách nhiệm rõ ràng. Tầng ngôn ngữ M1 thực hiện chuẩn hóa chuỗi,
tách từ theo so khớp cực đại trên từ điển Viet74K \cite{viet74k} và lọc từ dừng
tầng~1 kèm danh sách bảo vệ phủ định. Tầng truy xuất M2 biến dãy token thành chỉ
mục đảo, tính trọng số TF-IDF theo mô hình không gian vector \cite{salton1975vsm}
và xếp hạng bằng tương đồng cosin, đồng thời cài đặt độ đo Giá trị phân biệt từ
theo hướng tiếp cận của Salton và cộng sự \cite{salton1975theory} làm tầng lọc thứ
hai. Tầng tích hợp M3 cung cấp ứng dụng minh họa, hợp đồng API và bộ công cụ đóng
gói.

Bốn cấu hình thực nghiệm C1--C4 đã được dựng trên cùng một kho ngữ liệu gồm 2.500
bài thuộc 10 lĩnh vực và cùng một bộ 20 truy vấn: C1 không lọc, C2 lọc bằng danh
sách tĩnh, C3 lọc bằng hư từ tầng~1 của M1, và C4 lọc tầng~1 cộng thêm tầng~2 bằng
TDV. Các chỉ mục đã được dựng lại trên Maximum Matching chính thức của M1, với
metadata ghi rõ tokenizer, tài nguyên và hàm băm nguồn, nên so sánh giữa bốn cấu
hình là so sánh đúng một yếu tố: chiến lược lọc từ dừng.

\section{Những gì đã được kiểm chứng}

Phần tính toán cốt lõi đã được kiểm chứng bằng các đường tham chiếu độc lập, không
chỉ bằng các ca kiểm thử tự viết:

\begin{itemize}
  \item \textbf{TDV.} Đường tính nhanh trên cấu trúc thưa khớp với đường tham
        chiếu dựng ma trận cosin đầy đủ sau khi xóa từng chiều và tái chuẩn hóa,
        trong sai số số thực dưới $10^{-15}$. Bộ kiểm thử còn dựng sẵn một hàm mô
        phỏng ``cách làm sai'' (xóa chiều mà không tái chuẩn hóa) và khẳng định kết
        quả của engine khác kết quả sai đó, qua đó chứng minh tầm quan trọng của
        bước tái chuẩn hóa.
  \item \textbf{Chỉ số đánh giá.} Các định nghĩa $P@K$, $Recall@K$, $F1@K$, $AP$ và
        $NDCG@K$ được kiểm chứng bằng fixture tính tay; trọng số TF-IDF được đối
        soát với thư viện scikit-learn \cite{sklearn_tfidf} trên cùng cấu hình
        analyzer, sublinear TF, IDF làm mượt và chuẩn hóa L2.
  \item \textbf{Tách từ.} Ba phương pháp FMM/BMM/BiMM được kiểm chứng bằng từ điển
        fixture minh bạch, với quy tắc phân xử của BiMM được phát biểu tường minh.
        Vết thuật toán cho thấy ứng viên và quyết định tại từng bước; dãy token do
        \texttt{tokenize} trả về được kiểm riêng bằng fixture, không dùng các dòng
        ứng viên được đánh dấu của BMM để tái dựng kết quả.
  \item \textbf{Toàn vẹn dữ liệu.} Bộ đánh giá từ chối ghi kết quả chính thức khi
        hàm băm kho ngữ liệu, hàm băm tập truy vấn, trạng thái nhãn hoặc metadata
        tiền xử lý không đạt yêu cầu. Cơ chế khóa hàm băm tập truy vấn ngăn việc
        dùng nhầm nhãn cho một phiên bản truy vấn đã bị sửa.
  \item \textbf{Độ trễ và dung lượng.} Cả bốn cấu hình đã được đo độ trễ truy vấn và
        dung lượng chỉ mục trên máy cục bộ, với quy trình đo có bước khởi động làm
        nóng và có báo cáo theo từng truy vấn. Số đo chi tiết cùng trạng thái đạt
        hoặc không đạt mục tiêu được trình bày ở Chương~4.
\end{itemize}

\section{Những gì chưa đo được}

Báo cáo nêu rõ hai khoảng trống dữ liệu, và không thay thế chúng bằng kết luận định
tính:

\begin{enumerate}
  \item \textbf{Chỉ số chất lượng truy xuất chưa được đo.} Pool đánh giá có hàng
        trăm cặp ứng viên do ba đường cơ sở đóng góp, nhưng tại thời điểm viết
        \textbf{chưa có cặp nào được người gán xác nhận}; trạng thái nhãn là
        \texttt{pending\_human\_annotation}. Vì vậy bộ đánh giá ghi
        \texttt{metrics\_status = unmeasured} và cố ý không xuất bảng $P@K, R@K,
        F1@K, MAP, NDCG@K$. Không có kết luận nào về việc cấu hình nào cho chất
        lượng truy xuất tốt hơn, kể cả khi dung lượng chỉ mục của một cấu hình nhỏ
        hơn hay truy vấn của nó nhanh hơn.
  \item \textbf{Chất lượng tách từ chưa được đo.} Bộ mẫu 150 câu (15 câu mỗi lĩnh
        vực) rút từ kho ngữ liệu đã khóa vẫn ở trạng thái chờ người gán. Do đó
        không có chỉ số $F1$ tách từ nào cho cả ba phương pháp. Gợi ý của máy chỉ
        là công cụ hỗ trợ gán nhãn, không phải nhãn vàng.
\end{enumerate}

Việc giữ nguyên trạng thái ``chưa đo'' là một phần của kết quả: nó buộc mọi kết
luận về chất lượng phải chờ dữ liệu người, thay vì dựa vào các con số thời gian và
dung lượng vốn không phản ánh mức độ liên quan.

\section{Hạn chế}

Các hạn chế dưới đây cần được đọc kèm với các kết quả đã nêu:

\begin{itemize}
  \item \textbf{Phụ thuộc từ điển.} Chất lượng tách từ gắn chặt với từ điển đóng
        gói. Đổi từ điển thì mọi số đo tách từ phải được đo lại; từ điển rộng có
        thể chứa cả thành ngữ không đồng nghĩa với một phân đoạn từ tốt.
  \item \textbf{Lọc quá tay.} Cơ chế lọc theo tầng có thể loại cả những từ mang
        nghĩa với một số truy vấn cụ thể. Dung lượng nhỏ và truy vấn nhanh không
        chứng minh chất lượng tốt; vấn đề này chỉ giải quyết được khi có nhãn.
  \item \textbf{Phủ định ở mức túi từ.} Giữ lại từ phủ định tránh mất dấu hiệu phủ
        định ở mức từ vựng, nhưng mô hình túi từ vẫn không suy luận được ngữ nghĩa
        phủ định.
  \item \textbf{Quy mô đánh giá nhỏ và có thiên lệch.} Bộ truy vấn chỉ có 20 câu,
        khoảng hai câu mỗi lĩnh vực; pool được gộp từ chính các đường cơ sở nên có
        thể bỏ sót tài liệu liên quan không được đường nào trả về. Bộ mẫu tách từ
        150 câu, nếu chỉ một người gán, cũng không đo được đồng thuận giữa người gán.
  \item \textbf{Chi phí dự phòng.} Cơ chế dự phòng cho C4 khi mọi token truy vấn
        đều đã bị lọc đòi hỏi giữ thêm một chỉ mục C3. Việc giữ hai chỉ mục làm
        tổng dung lượng triển khai tăng lên, trái với mục tiêu giảm dung lượng của
        một chỉ mục đơn. Số liệu đo riêng chỉ mục C4 thuần và gói C4 cộng C3 được
        trình bày ở Chương~4.
  \item \textbf{Điều kiện đo độ trễ.} Độ trễ được đo trên tiến trình đã làm nóng
        trên một máy đơn, không phải cam kết về độ trễ trên mọi máy hay mọi kích
        thước kho ngữ liệu.
\end{itemize}

\section{Hướng phát triển}

Các bước tiếp theo được xác định theo thứ tự ưu tiên rõ ràng:

\begin{enumerate}
  \item Gán nhãn và xác nhận toàn bộ pool liên quan, sau đó chạy bộ đánh giá chính
        thức để có $P@K, R@K, F1@K, MAP, NDCG@K$ cho bốn cấu hình và cho từng lĩnh
        vực.
  \item Hoàn tất gán nhãn 150 câu mẫu tách từ, rồi chạy bộ đo khớp span để có $F1$
        cho FMM, BMM, BiMM.
  \item So sánh thêm với một bộ tách từ ngoài đã được công bố, trong điều kiện cho
        phép, để đặt kết quả của nhóm vào một mốc tham chiếu.
  \item Mở rộng bộ mẫu và số người gán để đo đồng thuận giữa người gán và để giảm
        độ dao động của chỉ số theo lĩnh vực.
  \item Nghiên cứu thay thế túi từ bằng một biểu diễn xử lý được phủ định ở mức
        ngữ nghĩa, và đánh giá lại tác động của việc bảo vệ guard.
\end{enumerate}

Cho tới khi các bước gán nhãn hoàn tất, kết luận đúng mực của đồ án là: hệ thống đã
hoàn tất về mặt triển khai và đã kiểm chứng được phần tính toán, thời gian và dung
lượng; phần kết luận về chất lượng truy xuất và chất lượng tách từ vẫn đang chờ
nhãn người.
. Chỉ dùng các gói đã khai báo.
\chapter{KẾT LUẬN VÀ HƯỚNG PHÁT TRIỂN}
\label{ch:ch5-ketluan}

\section{Kết quả đạt được}

Đồ án đã hoàn thành một hệ thống truy xuất văn bản tiếng Việt chạy hoàn toàn ngoại
tuyến, gồm ba tầng trách nhiệm rõ ràng. Tầng ngôn ngữ M1 thực hiện chuẩn hóa chuỗi,
tách từ theo so khớp cực đại trên từ điển Viet74K \cite{viet74k} và lọc từ dừng
tầng~1 kèm danh sách bảo vệ phủ định. Tầng truy xuất M2 biến dãy token thành chỉ
mục đảo, tính trọng số TF-IDF theo mô hình không gian vector \cite{salton1975vsm}
và xếp hạng bằng tương đồng cosin, đồng thời cài đặt độ đo Giá trị phân biệt từ
theo hướng tiếp cận của Salton và cộng sự \cite{salton1975theory} làm tầng lọc thứ
hai. Tầng tích hợp M3 cung cấp ứng dụng minh họa, hợp đồng API và bộ công cụ đóng
gói.

Bốn cấu hình thực nghiệm C1--C4 đã được dựng trên cùng một kho ngữ liệu gồm 2.500
bài thuộc 10 lĩnh vực và cùng một bộ 20 truy vấn: C1 không lọc, C2 lọc bằng danh
sách tĩnh, C3 lọc bằng hư từ tầng~1 của M1, và C4 lọc tầng~1 cộng thêm tầng~2 bằng
TDV. Các chỉ mục đã được dựng lại trên Maximum Matching chính thức của M1, với
metadata ghi rõ tokenizer, tài nguyên và hàm băm nguồn, nên so sánh giữa bốn cấu
hình là so sánh đúng một yếu tố: chiến lược lọc từ dừng.

\section{Những gì đã được kiểm chứng}

Phần tính toán cốt lõi đã được kiểm chứng bằng các đường tham chiếu độc lập, không
chỉ bằng các ca kiểm thử tự viết:

\begin{itemize}
  \item \textbf{TDV.} Đường tính nhanh trên cấu trúc thưa khớp với đường tham
        chiếu dựng ma trận cosin đầy đủ sau khi xóa từng chiều và tái chuẩn hóa,
        trong sai số số thực dưới $10^{-15}$. Bộ kiểm thử còn dựng sẵn một hàm mô
        phỏng ``cách làm sai'' (xóa chiều mà không tái chuẩn hóa) và khẳng định kết
        quả của engine khác kết quả sai đó, qua đó chứng minh tầm quan trọng của
        bước tái chuẩn hóa.
  \item \textbf{Chỉ số đánh giá.} Các định nghĩa $P@K$, $Recall@K$, $F1@K$, $AP$ và
        $NDCG@K$ được kiểm chứng bằng fixture tính tay; trọng số TF-IDF được đối
        soát với thư viện scikit-learn \cite{sklearn_tfidf} trên cùng cấu hình
        analyzer, sublinear TF, IDF làm mượt và chuẩn hóa L2.
  \item \textbf{Tách từ.} Ba phương pháp FMM/BMM/BiMM được kiểm chứng bằng từ điển
        fixture minh bạch, với quy tắc phân xử của BiMM được phát biểu tường minh.
        Vết thuật toán cho thấy ứng viên và quyết định tại từng bước; dãy token do
        \texttt{tokenize} trả về được kiểm riêng bằng fixture, không dùng các dòng
        ứng viên được đánh dấu của BMM để tái dựng kết quả.
  \item \textbf{Toàn vẹn dữ liệu.} Bộ đánh giá từ chối ghi kết quả chính thức khi
        hàm băm kho ngữ liệu, hàm băm tập truy vấn, trạng thái nhãn hoặc metadata
        tiền xử lý không đạt yêu cầu. Cơ chế khóa hàm băm tập truy vấn ngăn việc
        dùng nhầm nhãn cho một phiên bản truy vấn đã bị sửa.
  \item \textbf{Độ trễ và dung lượng.} Cả bốn cấu hình đã được đo độ trễ truy vấn và
        dung lượng chỉ mục trên máy cục bộ, với quy trình đo có bước khởi động làm
        nóng và có báo cáo theo từng truy vấn. Số đo chi tiết cùng trạng thái đạt
        hoặc không đạt mục tiêu được trình bày ở Chương~4.
\end{itemize}

\section{Những gì chưa đo được}

Báo cáo nêu rõ hai khoảng trống dữ liệu, và không thay thế chúng bằng kết luận định
tính:

\begin{enumerate}
  \item \textbf{Chỉ số chất lượng truy xuất chưa được đo.} Pool đánh giá có hàng
        trăm cặp ứng viên do ba đường cơ sở đóng góp, nhưng tại thời điểm viết
        \textbf{chưa có cặp nào được người gán xác nhận}; trạng thái nhãn là
        \texttt{pending\_human\_annotation}. Vì vậy bộ đánh giá ghi
        \texttt{metrics\_status = unmeasured} và cố ý không xuất bảng $P@K, R@K,
        F1@K, MAP, NDCG@K$. Không có kết luận nào về việc cấu hình nào cho chất
        lượng truy xuất tốt hơn, kể cả khi dung lượng chỉ mục của một cấu hình nhỏ
        hơn hay truy vấn của nó nhanh hơn.
  \item \textbf{Chất lượng tách từ chưa được đo.} Bộ mẫu 150 câu (15 câu mỗi lĩnh
        vực) rút từ kho ngữ liệu đã khóa vẫn ở trạng thái chờ người gán. Do đó
        không có chỉ số $F1$ tách từ nào cho cả ba phương pháp. Gợi ý của máy chỉ
        là công cụ hỗ trợ gán nhãn, không phải nhãn vàng.
\end{enumerate}

Việc giữ nguyên trạng thái ``chưa đo'' là một phần của kết quả: nó buộc mọi kết
luận về chất lượng phải chờ dữ liệu người, thay vì dựa vào các con số thời gian và
dung lượng vốn không phản ánh mức độ liên quan.

\section{Hạn chế}

Các hạn chế dưới đây cần được đọc kèm với các kết quả đã nêu:

\begin{itemize}
  \item \textbf{Phụ thuộc từ điển.} Chất lượng tách từ gắn chặt với từ điển đóng
        gói. Đổi từ điển thì mọi số đo tách từ phải được đo lại; từ điển rộng có
        thể chứa cả thành ngữ không đồng nghĩa với một phân đoạn từ tốt.
  \item \textbf{Lọc quá tay.} Cơ chế lọc theo tầng có thể loại cả những từ mang
        nghĩa với một số truy vấn cụ thể. Dung lượng nhỏ và truy vấn nhanh không
        chứng minh chất lượng tốt; vấn đề này chỉ giải quyết được khi có nhãn.
  \item \textbf{Phủ định ở mức túi từ.} Giữ lại từ phủ định tránh mất dấu hiệu phủ
        định ở mức từ vựng, nhưng mô hình túi từ vẫn không suy luận được ngữ nghĩa
        phủ định.
  \item \textbf{Quy mô đánh giá nhỏ và có thiên lệch.} Bộ truy vấn chỉ có 20 câu,
        khoảng hai câu mỗi lĩnh vực; pool được gộp từ chính các đường cơ sở nên có
        thể bỏ sót tài liệu liên quan không được đường nào trả về. Bộ mẫu tách từ
        150 câu, nếu chỉ một người gán, cũng không đo được đồng thuận giữa người gán.
  \item \textbf{Chi phí dự phòng.} Cơ chế dự phòng cho C4 khi mọi token truy vấn
        đều đã bị lọc đòi hỏi giữ thêm một chỉ mục C3. Việc giữ hai chỉ mục làm
        tổng dung lượng triển khai tăng lên, trái với mục tiêu giảm dung lượng của
        một chỉ mục đơn. Số liệu đo riêng chỉ mục C4 thuần và gói C4 cộng C3 được
        trình bày ở Chương~4.
  \item \textbf{Điều kiện đo độ trễ.} Độ trễ được đo trên tiến trình đã làm nóng
        trên một máy đơn, không phải cam kết về độ trễ trên mọi máy hay mọi kích
        thước kho ngữ liệu.
\end{itemize}

\section{Hướng phát triển}

Các bước tiếp theo được xác định theo thứ tự ưu tiên rõ ràng:

\begin{enumerate}
  \item Gán nhãn và xác nhận toàn bộ pool liên quan, sau đó chạy bộ đánh giá chính
        thức để có $P@K, R@K, F1@K, MAP, NDCG@K$ cho bốn cấu hình và cho từng lĩnh
        vực.
  \item Hoàn tất gán nhãn 150 câu mẫu tách từ, rồi chạy bộ đo khớp span để có $F1$
        cho FMM, BMM, BiMM.
  \item So sánh thêm với một bộ tách từ ngoài đã được công bố, trong điều kiện cho
        phép, để đặt kết quả của nhóm vào một mốc tham chiếu.
  \item Mở rộng bộ mẫu và số người gán để đo đồng thuận giữa người gán và để giảm
        độ dao động của chỉ số theo lĩnh vực.
  \item Nghiên cứu thay thế túi từ bằng một biểu diễn xử lý được phủ định ở mức
        ngữ nghĩa, và đánh giá lại tác động của việc bảo vệ guard.
\end{enumerate}

Cho tới khi các bước gán nhãn hoàn tất, kết luận đúng mực của đồ án là: hệ thống đã
hoàn tất về mặt triển khai và đã kiểm chứng được phần tính toán, thời gian và dung
lượng; phần kết luận về chất lượng truy xuất và chất lượng tách từ vẫn đang chờ
nhãn người.
. Chỉ dùng các gói đã khai báo.
\chapter{KẾT LUẬN VÀ HƯỚNG PHÁT TRIỂN}
\label{ch:ch5-ketluan}

\section{Kết quả đạt được}

Đồ án đã hoàn thành một hệ thống truy xuất văn bản tiếng Việt chạy hoàn toàn ngoại
tuyến, gồm ba tầng trách nhiệm rõ ràng. Tầng ngôn ngữ M1 thực hiện chuẩn hóa chuỗi,
tách từ theo so khớp cực đại trên từ điển Viet74K \cite{viet74k} và lọc từ dừng
tầng~1 kèm danh sách bảo vệ phủ định. Tầng truy xuất M2 biến dãy token thành chỉ
mục đảo, tính trọng số TF-IDF theo mô hình không gian vector \cite{salton1975vsm}
và xếp hạng bằng tương đồng cosin, đồng thời cài đặt độ đo Giá trị phân biệt từ
theo hướng tiếp cận của Salton và cộng sự \cite{salton1975theory} làm tầng lọc thứ
hai. Tầng tích hợp M3 cung cấp ứng dụng minh họa, hợp đồng API và bộ công cụ đóng
gói.

Bốn cấu hình thực nghiệm C1--C4 đã được dựng trên cùng một kho ngữ liệu gồm 2.500
bài thuộc 10 lĩnh vực và cùng một bộ 20 truy vấn: C1 không lọc, C2 lọc bằng danh
sách tĩnh, C3 lọc bằng hư từ tầng~1 của M1, và C4 lọc tầng~1 cộng thêm tầng~2 bằng
TDV. Các chỉ mục đã được dựng lại trên Maximum Matching chính thức của M1, với
metadata ghi rõ tokenizer, tài nguyên và hàm băm nguồn, nên so sánh giữa bốn cấu
hình là so sánh đúng một yếu tố: chiến lược lọc từ dừng.

\section{Những gì đã được kiểm chứng}

Phần tính toán cốt lõi đã được kiểm chứng bằng các đường tham chiếu độc lập, không
chỉ bằng các ca kiểm thử tự viết:

\begin{itemize}
  \item \textbf{TDV.} Đường tính nhanh trên cấu trúc thưa khớp với đường tham
        chiếu dựng ma trận cosin đầy đủ sau khi xóa từng chiều và tái chuẩn hóa,
        trong sai số số thực dưới $10^{-15}$. Bộ kiểm thử còn dựng sẵn một hàm mô
        phỏng ``cách làm sai'' (xóa chiều mà không tái chuẩn hóa) và khẳng định kết
        quả của engine khác kết quả sai đó, qua đó chứng minh tầm quan trọng của
        bước tái chuẩn hóa.
  \item \textbf{Chỉ số đánh giá.} Các định nghĩa $P@K$, $Recall@K$, $F1@K$, $AP$ và
        $NDCG@K$ được kiểm chứng bằng fixture tính tay; trọng số TF-IDF được đối
        soát với thư viện scikit-learn \cite{sklearn_tfidf} trên cùng cấu hình
        analyzer, sublinear TF, IDF làm mượt và chuẩn hóa L2.
  \item \textbf{Tách từ.} Ba phương pháp FMM/BMM/BiMM được kiểm chứng bằng từ điển
        fixture minh bạch, với quy tắc phân xử của BiMM được phát biểu tường minh.
        Vết thuật toán cho thấy ứng viên và quyết định tại từng bước; dãy token do
        \texttt{tokenize} trả về được kiểm riêng bằng fixture, không dùng các dòng
        ứng viên được đánh dấu của BMM để tái dựng kết quả.
  \item \textbf{Toàn vẹn dữ liệu.} Bộ đánh giá từ chối ghi kết quả chính thức khi
        hàm băm kho ngữ liệu, hàm băm tập truy vấn, trạng thái nhãn hoặc metadata
        tiền xử lý không đạt yêu cầu. Cơ chế khóa hàm băm tập truy vấn ngăn việc
        dùng nhầm nhãn cho một phiên bản truy vấn đã bị sửa.
  \item \textbf{Độ trễ và dung lượng.} Cả bốn cấu hình đã được đo độ trễ truy vấn và
        dung lượng chỉ mục trên máy cục bộ, với quy trình đo có bước khởi động làm
        nóng và có báo cáo theo từng truy vấn. Số đo chi tiết cùng trạng thái đạt
        hoặc không đạt mục tiêu được trình bày ở Chương~4.
\end{itemize}

\section{Những gì chưa đo được}

Báo cáo nêu rõ hai khoảng trống dữ liệu, và không thay thế chúng bằng kết luận định
tính:

\begin{enumerate}
  \item \textbf{Chỉ số chất lượng truy xuất chưa được đo.} Pool đánh giá có hàng
        trăm cặp ứng viên do ba đường cơ sở đóng góp, nhưng tại thời điểm viết
        \textbf{chưa có cặp nào được người gán xác nhận}; trạng thái nhãn là
        \texttt{pending\_human\_annotation}. Vì vậy bộ đánh giá ghi
        \texttt{metrics\_status = unmeasured} và cố ý không xuất bảng $P@K, R@K,
        F1@K, MAP, NDCG@K$. Không có kết luận nào về việc cấu hình nào cho chất
        lượng truy xuất tốt hơn, kể cả khi dung lượng chỉ mục của một cấu hình nhỏ
        hơn hay truy vấn của nó nhanh hơn.
  \item \textbf{Chất lượng tách từ chưa được đo.} Bộ mẫu 150 câu (15 câu mỗi lĩnh
        vực) rút từ kho ngữ liệu đã khóa vẫn ở trạng thái chờ người gán. Do đó
        không có chỉ số $F1$ tách từ nào cho cả ba phương pháp. Gợi ý của máy chỉ
        là công cụ hỗ trợ gán nhãn, không phải nhãn vàng.
\end{enumerate}

Việc giữ nguyên trạng thái ``chưa đo'' là một phần của kết quả: nó buộc mọi kết
luận về chất lượng phải chờ dữ liệu người, thay vì dựa vào các con số thời gian và
dung lượng vốn không phản ánh mức độ liên quan.

\section{Hạn chế}

Các hạn chế dưới đây cần được đọc kèm với các kết quả đã nêu:

\begin{itemize}
  \item \textbf{Phụ thuộc từ điển.} Chất lượng tách từ gắn chặt với từ điển đóng
        gói. Đổi từ điển thì mọi số đo tách từ phải được đo lại; từ điển rộng có
        thể chứa cả thành ngữ không đồng nghĩa với một phân đoạn từ tốt.
  \item \textbf{Lọc quá tay.} Cơ chế lọc theo tầng có thể loại cả những từ mang
        nghĩa với một số truy vấn cụ thể. Dung lượng nhỏ và truy vấn nhanh không
        chứng minh chất lượng tốt; vấn đề này chỉ giải quyết được khi có nhãn.
  \item \textbf{Phủ định ở mức túi từ.} Giữ lại từ phủ định tránh mất dấu hiệu phủ
        định ở mức từ vựng, nhưng mô hình túi từ vẫn không suy luận được ngữ nghĩa
        phủ định.
  \item \textbf{Quy mô đánh giá nhỏ và có thiên lệch.} Bộ truy vấn chỉ có 20 câu,
        khoảng hai câu mỗi lĩnh vực; pool được gộp từ chính các đường cơ sở nên có
        thể bỏ sót tài liệu liên quan không được đường nào trả về. Bộ mẫu tách từ
        150 câu, nếu chỉ một người gán, cũng không đo được đồng thuận giữa người gán.
  \item \textbf{Chi phí dự phòng.} Cơ chế dự phòng cho C4 khi mọi token truy vấn
        đều đã bị lọc đòi hỏi giữ thêm một chỉ mục C3. Việc giữ hai chỉ mục làm
        tổng dung lượng triển khai tăng lên, trái với mục tiêu giảm dung lượng của
        một chỉ mục đơn. Số liệu đo riêng chỉ mục C4 thuần và gói C4 cộng C3 được
        trình bày ở Chương~4.
  \item \textbf{Điều kiện đo độ trễ.} Độ trễ được đo trên tiến trình đã làm nóng
        trên một máy đơn, không phải cam kết về độ trễ trên mọi máy hay mọi kích
        thước kho ngữ liệu.
\end{itemize}

\section{Hướng phát triển}

Các bước tiếp theo được xác định theo thứ tự ưu tiên rõ ràng:

\begin{enumerate}
  \item Gán nhãn và xác nhận toàn bộ pool liên quan, sau đó chạy bộ đánh giá chính
        thức để có $P@K, R@K, F1@K, MAP, NDCG@K$ cho bốn cấu hình và cho từng lĩnh
        vực.
  \item Hoàn tất gán nhãn 150 câu mẫu tách từ, rồi chạy bộ đo khớp span để có $F1$
        cho FMM, BMM, BiMM.
  \item So sánh thêm với một bộ tách từ ngoài đã được công bố, trong điều kiện cho
        phép, để đặt kết quả của nhóm vào một mốc tham chiếu.
  \item Mở rộng bộ mẫu và số người gán để đo đồng thuận giữa người gán và để giảm
        độ dao động của chỉ số theo lĩnh vực.
  \item Nghiên cứu thay thế túi từ bằng một biểu diễn xử lý được phủ định ở mức
        ngữ nghĩa, và đánh giá lại tác động của việc bảo vệ guard.
\end{enumerate}

Cho tới khi các bước gán nhãn hoàn tất, kết luận đúng mực của đồ án là: hệ thống đã
hoàn tất về mặt triển khai và đã kiểm chứng được phần tính toán, thời gian và dung
lượng; phần kết luận về chất lượng truy xuất và chất lượng tách từ vẫn đang chờ
nhãn người.
. Chỉ dùng các gói đã khai báo.
\chapter{KẾT LUẬN VÀ HƯỚNG PHÁT TRIỂN}
\label{ch:ch5-ketluan}

\section{Kết quả đạt được}

Đồ án đã hoàn thành một hệ thống truy xuất văn bản tiếng Việt chạy hoàn toàn ngoại
tuyến, gồm ba tầng trách nhiệm rõ ràng. Tầng ngôn ngữ M1 thực hiện chuẩn hóa chuỗi,
tách từ theo so khớp cực đại trên từ điển Viet74K \cite{viet74k} và lọc từ dừng
tầng~1 kèm danh sách bảo vệ phủ định. Tầng truy xuất M2 biến dãy token thành chỉ
mục đảo, tính trọng số TF-IDF theo mô hình không gian vector \cite{salton1975vsm}
và xếp hạng bằng tương đồng cosin, đồng thời cài đặt độ đo Giá trị phân biệt từ
theo hướng tiếp cận của Salton và cộng sự \cite{salton1975theory} làm tầng lọc thứ
hai. Tầng tích hợp M3 cung cấp ứng dụng minh họa, hợp đồng API và bộ công cụ đóng
gói.

Bốn cấu hình thực nghiệm C1--C4 đã được dựng trên cùng một kho ngữ liệu gồm 2.500
bài thuộc 10 lĩnh vực và cùng một bộ 20 truy vấn: C1 không lọc, C2 lọc bằng danh
sách tĩnh, C3 lọc bằng hư từ tầng~1 của M1, và C4 lọc tầng~1 cộng thêm tầng~2 bằng
TDV. Các chỉ mục đã được dựng lại trên Maximum Matching chính thức của M1, với
metadata ghi rõ tokenizer, tài nguyên và hàm băm nguồn, nên so sánh giữa bốn cấu
hình là so sánh đúng một yếu tố: chiến lược lọc từ dừng.

\section{Những gì đã được kiểm chứng}

Phần tính toán cốt lõi đã được kiểm chứng bằng các đường tham chiếu độc lập, không
chỉ bằng các ca kiểm thử tự viết:

\begin{itemize}
  \item \textbf{TDV.} Đường tính nhanh trên cấu trúc thưa khớp với đường tham
        chiếu dựng ma trận cosin đầy đủ sau khi xóa từng chiều và tái chuẩn hóa,
        trong sai số số thực dưới $10^{-15}$. Bộ kiểm thử còn dựng sẵn một hàm mô
        phỏng ``cách làm sai'' (xóa chiều mà không tái chuẩn hóa) và khẳng định kết
        quả của engine khác kết quả sai đó, qua đó chứng minh tầm quan trọng của
        bước tái chuẩn hóa.
  \item \textbf{Chỉ số đánh giá.} Các định nghĩa $P@K$, $Recall@K$, $F1@K$, $AP$ và
        $NDCG@K$ được kiểm chứng bằng fixture tính tay; trọng số TF-IDF được đối
        soát với thư viện scikit-learn \cite{sklearn_tfidf} trên cùng cấu hình
        analyzer, sublinear TF, IDF làm mượt và chuẩn hóa L2.
  \item \textbf{Tách từ.} Ba phương pháp FMM/BMM/BiMM được kiểm chứng bằng từ điển
        fixture minh bạch, với quy tắc phân xử của BiMM được phát biểu tường minh.
        Vết thuật toán cho thấy ứng viên và quyết định tại từng bước; dãy token do
        \texttt{tokenize} trả về được kiểm riêng bằng fixture, không dùng các dòng
        ứng viên được đánh dấu của BMM để tái dựng kết quả.
  \item \textbf{Toàn vẹn dữ liệu.} Bộ đánh giá từ chối ghi kết quả chính thức khi
        hàm băm kho ngữ liệu, hàm băm tập truy vấn, trạng thái nhãn hoặc metadata
        tiền xử lý không đạt yêu cầu. Cơ chế khóa hàm băm tập truy vấn ngăn việc
        dùng nhầm nhãn cho một phiên bản truy vấn đã bị sửa.
  \item \textbf{Độ trễ và dung lượng.} Cả bốn cấu hình đã được đo độ trễ truy vấn và
        dung lượng chỉ mục trên máy cục bộ, với quy trình đo có bước khởi động làm
        nóng và có báo cáo theo từng truy vấn. Số đo chi tiết cùng trạng thái đạt
        hoặc không đạt mục tiêu được trình bày ở Chương~4.
\end{itemize}

\section{Những gì chưa đo được}

Báo cáo nêu rõ hai khoảng trống dữ liệu, và không thay thế chúng bằng kết luận định
tính:

\begin{enumerate}
  \item \textbf{Chỉ số chất lượng truy xuất chưa được đo.} Pool đánh giá có hàng
        trăm cặp ứng viên do ba đường cơ sở đóng góp, nhưng tại thời điểm viết
        \textbf{chưa có cặp nào được người gán xác nhận}; trạng thái nhãn là
        \texttt{pending\_human\_annotation}. Vì vậy bộ đánh giá ghi
        \texttt{metrics\_status = unmeasured} và cố ý không xuất bảng $P@K, R@K,
        F1@K, MAP, NDCG@K$. Không có kết luận nào về việc cấu hình nào cho chất
        lượng truy xuất tốt hơn, kể cả khi dung lượng chỉ mục của một cấu hình nhỏ
        hơn hay truy vấn của nó nhanh hơn.
  \item \textbf{Chất lượng tách từ chưa được đo.} Bộ mẫu 150 câu (15 câu mỗi lĩnh
        vực) rút từ kho ngữ liệu đã khóa vẫn ở trạng thái chờ người gán. Do đó
        không có chỉ số $F1$ tách từ nào cho cả ba phương pháp. Gợi ý của máy chỉ
        là công cụ hỗ trợ gán nhãn, không phải nhãn vàng.
\end{enumerate}

Việc giữ nguyên trạng thái ``chưa đo'' là một phần của kết quả: nó buộc mọi kết
luận về chất lượng phải chờ dữ liệu người, thay vì dựa vào các con số thời gian và
dung lượng vốn không phản ánh mức độ liên quan.

\section{Hạn chế}

Các hạn chế dưới đây cần được đọc kèm với các kết quả đã nêu:

\begin{itemize}
  \item \textbf{Phụ thuộc từ điển.} Chất lượng tách từ gắn chặt với từ điển đóng
        gói. Đổi từ điển thì mọi số đo tách từ phải được đo lại; từ điển rộng có
        thể chứa cả thành ngữ không đồng nghĩa với một phân đoạn từ tốt.
  \item \textbf{Lọc quá tay.} Cơ chế lọc theo tầng có thể loại cả những từ mang
        nghĩa với một số truy vấn cụ thể. Dung lượng nhỏ và truy vấn nhanh không
        chứng minh chất lượng tốt; vấn đề này chỉ giải quyết được khi có nhãn.
  \item \textbf{Phủ định ở mức túi từ.} Giữ lại từ phủ định tránh mất dấu hiệu phủ
        định ở mức từ vựng, nhưng mô hình túi từ vẫn không suy luận được ngữ nghĩa
        phủ định.
  \item \textbf{Quy mô đánh giá nhỏ và có thiên lệch.} Bộ truy vấn chỉ có 20 câu,
        khoảng hai câu mỗi lĩnh vực; pool được gộp từ chính các đường cơ sở nên có
        thể bỏ sót tài liệu liên quan không được đường nào trả về. Bộ mẫu tách từ
        150 câu, nếu chỉ một người gán, cũng không đo được đồng thuận giữa người gán.
  \item \textbf{Chi phí dự phòng.} Cơ chế dự phòng cho C4 khi mọi token truy vấn
        đều đã bị lọc đòi hỏi giữ thêm một chỉ mục C3. Việc giữ hai chỉ mục làm
        tổng dung lượng triển khai tăng lên, trái với mục tiêu giảm dung lượng của
        một chỉ mục đơn. Số liệu đo riêng chỉ mục C4 thuần và gói C4 cộng C3 được
        trình bày ở Chương~4.
  \item \textbf{Điều kiện đo độ trễ.} Độ trễ được đo trên tiến trình đã làm nóng
        trên một máy đơn, không phải cam kết về độ trễ trên mọi máy hay mọi kích
        thước kho ngữ liệu.
\end{itemize}

\section{Hướng phát triển}

Các bước tiếp theo được xác định theo thứ tự ưu tiên rõ ràng:

\begin{enumerate}
  \item Gán nhãn và xác nhận toàn bộ pool liên quan, sau đó chạy bộ đánh giá chính
        thức để có $P@K, R@K, F1@K, MAP, NDCG@K$ cho bốn cấu hình và cho từng lĩnh
        vực.
  \item Hoàn tất gán nhãn 150 câu mẫu tách từ, rồi chạy bộ đo khớp span để có $F1$
        cho FMM, BMM, BiMM.
  \item So sánh thêm với một bộ tách từ ngoài đã được công bố, trong điều kiện cho
        phép, để đặt kết quả của nhóm vào một mốc tham chiếu.
  \item Mở rộng bộ mẫu và số người gán để đo đồng thuận giữa người gán và để giảm
        độ dao động của chỉ số theo lĩnh vực.
  \item Nghiên cứu thay thế túi từ bằng một biểu diễn xử lý được phủ định ở mức
        ngữ nghĩa, và đánh giá lại tác động của việc bảo vệ guard.
\end{enumerate}

Cho tới khi các bước gán nhãn hoàn tất, kết luận đúng mực của đồ án là: hệ thống đã
hoàn tất về mặt triển khai và đã kiểm chứng được phần tính toán, thời gian và dung
lượng; phần kết luận về chất lượng truy xuất và chất lượng tách từ vẫn đang chờ
nhãn người.
